\section{Method}
\label{sec:method}

\subsection{Overview}

Problem~\eqref{eq:compression_aware_problem} couples two objectives that live in
different spaces: bitrate is measured in the latent domain of a frozen neural
codec, while fidelity is measured in image space relative to the pretrained SR
reference \(x_{\mathrm{ref}}\). Direct joint optimisation is difficult because
both the SR network and the codec are highly nonlinear, and because the decision
variables are network parameters \(\theta\) whereas the constraints are expressed
on images.

We therefore use an ADMM-inspired variable-splitting strategy~\cite{boyd2011admm}.
The PSNR constraint on the SR image is enforced through an auxiliary image that
is projected onto a feasible quality ball; the PSNR constraint on the compressed
reconstruction is handled as a soft penalty; and the classical quadratic
coupling term is replaced in practice by a PSNR-domain surrogate so that all
loss terms operate on a comparable scale.

\subsection{Constraint reformulation}

Reducing bitrate tends to favour smoother, more compressible reconstructions,
while the fidelity constraints require preserving the high-frequency content of
\(x_{\mathrm{ref}}\). To separate these competing goals, introduce an auxiliary
image \(z\) constrained to the feasible set
\begin{equation}
  \mathcal{B}
  = \bigl\{ z : \mathrm{PSNR}(z, x_{\mathrm{ref}}) \ge P_{\mathrm{target}} \bigr\}.
\end{equation}
Geometrically, \(\mathcal{B}\) is a closed ball in image space centred at
\(x_{\mathrm{ref}}\) with radius \(\delta\) determined by \(P_{\mathrm{target}}\):
\begin{equation}
  \delta^2
  = \frac{C_{\mathrm{ch}} H W}{10^{P_{\mathrm{target}} / 10}},
  \label{eq:psnr_radius}
\end{equation}
where \(C_{\mathrm{ch}}\), \(H\), and \(W\) are the channel and spatial dimensions
of the SR image (pixel values in \([0,1]\)).

The compressed-image constraint is harder to enforce exactly because it depends
on the nonlinear map \(\mathrm{Dec}\circ\mathrm{Enc}\). We fold it into the
objective as a soft penalty. Define
\begin{equation}
  f(x) = R\!\bigl(\mathrm{Enc}(x)\bigr) + \lambda\,\phi(\hat{x}),
  \qquad
  \hat{x} = \mathrm{Dec}\!\bigl(\mathrm{Enc}(x)\bigr),
\end{equation}
with a dead-zone hinge around the target quality,
\begin{equation}
  \phi(\hat{x})
  = \max\!\left(
      \bigl|\mathrm{PSNR}(\hat{x}, x_{\mathrm{ref}}) - P_{\mathrm{target}}\bigr|
      - 0.5,\; 0
    \right).
  \label{eq:soft_penalty}
\end{equation}
The \(0.5\,\mathrm{dB}\) margin avoids unnecessary gradient updates when the
compressed reconstruction already meets the quality budget.

The reformulated problem is then
\begin{equation}
  \min_{x,\,z} \; f(x) + I_{\mathcal{B}}(z)
  \quad \text{subject to} \quad x = z,
  \label{eq:split_problem}
\end{equation}
where \(I_{\mathcal{B}}\) is the indicator of \(\mathcal{B}\)
(\(0\) on \(\mathcal{B}\), \(+\infty\) outside). Here \(x\) is responsible for
bitrate reduction and \(z\) for fidelity; the equality \(x=z\) couples the two.

\subsection{ADMM updates}

With a scaled dual variable \(t\) and penalty parameter \(\gamma>0\), an
augmented Lagrangian for~\eqref{eq:split_problem} is
\begin{equation}
  \mathcal{L}_{\gamma}(x, z, t)
  = f(x) + I_{\mathcal{B}}(z)
    + \frac{\gamma}{2}
      \Bigl(\|x - z - t\|_2^2 - \|t\|_2^2\Bigr).
\end{equation}
Alternating minimisation yields the ADMM-style updates
\begin{align}
  x^{k+1}
  &= \arg\min_{x} \;
     f(x) + \tfrac{\gamma}{2}\|x - (z^{k} + t^{k})\|_2^2,
     \label{eq:x_update} \\
  z^{k+1}
  &= \Pi_{\mathcal{B}}(x^{k+1} - t^{k}),
     \label{eq:z_update} \\
  t^{k+1}
  &= t^{k} + z^{k+1} - x^{k+1}.
     \label{eq:t_update}
\end{align}
In our setting the \(x\)-update is realised by gradient steps on the SR
parameters \(\theta\) (with \(x = f_{\theta}(x_{\mathrm{LR}})\)), while the codec
\(\mathrm{Enc},\mathrm{Dec}\) remains frozen. The \(z\)-update projects onto the
PSNR ball, and the dual update accumulates the residual \(z-x\).

\subsection{Projection onto the PSNR ball}

The \(z\)-step is a Euclidean projection onto \(\mathcal{B}\). Let
\begin{equation}
  d^{k} = x^{k+1} - t^{k},
  \qquad
  e^{k} = d^{k} - x_{\mathrm{ref}}.
\end{equation}
If \(\|e^{k}\|_2 \le \delta\), set \(z^{k+1} = d^{k}\). Otherwise project onto
the sphere boundary:
\begin{equation}
  z^{k+1}
  = x_{\mathrm{ref}} + \delta\,\frac{e^{k}}{\|e^{k}\|_2 + \varepsilon},
  \label{eq:projection_step}
\end{equation}
with a small \(\varepsilon>0\) for numerical stability. This closed form keeps
\(z^{k+1}\in\mathcal{B}\) at every outer iteration.

\subsection{Practical training loss}

A literal quadratic coupling \(\|x-(z+t)\|_2^2\) mixes poorly with bitrate
\(R(y)\) measured in bits per pixel: the two quantities differ by orders of
magnitude and make \(\gamma\) brittle. We therefore replace the quadratic term
by a PSNR-based surrogate on the consensus target \(v^{k}=z^{k}+t^{k}\). The
implemented inner-loop loss is
\begin{equation}
  \begin{aligned}
    \mathcal{L}
    &= 10\log_{10}\!\bigl(R(y)+\varepsilon\bigr) \\
    &\quad
      + \lambda_1\,\max\!\left(
          \bigl|\mathrm{PSNR}(\hat{x}, x_{\mathrm{ref}}) - P_{\mathrm{target}}\bigr|
          - 0.5,\; 0
        \right) \\
    &\quad
      + \lambda_2\,\max\!\left(
          \bigl|\mathrm{PSNR}(x, v^{k}) - P_{\mathrm{target}}\bigr|
          - 0.5,\; 0
        \right),
  \end{aligned}
  \label{eq:training_loss}
\end{equation}
where \(y=\mathrm{Enc}(x)\) and \(\hat{x}=\mathrm{Dec}(y)\). The first term
drives bitrate down on a log scale; the second softly enforces compressed-image
fidelity to \(x_{\mathrm{ref}}\); the third pulls the SR output toward the ADMM
consensus point in PSNR units. The outer \(z\)- and \(t\)-updates remain as
in~\eqref{eq:z_update}--\eqref{eq:t_update}.

We call the scheme ADMM-\emph{inspired}: it borrows variable splitting and
projection for controllability, without claiming formal convex ADMM convergence
for the deep nonlinear composite map.

\subsection{Per-image algorithm}

\begin{enumerate}
    \item Construct \(x_{\mathrm{LR}}\) (bicubic downsampling of an HR crop) and
    set \(x_{\mathrm{ref}}=f_{\theta_0}(x_{\mathrm{LR}})\). Freeze
    \(\mathrm{Enc},\mathrm{Dec}\).
    \item Choose trainable parameters: full NinaSR, full Swin2SR, or LoRA
    adapters on Swin2SR. Initialise \(z^{0}=x_{\mathrm{ref}}\), \(t^{0}=0\).
    \item For outer iteration \(k=0,\ldots,K-1\):
    \begin{enumerate}
        \item Set \(v^{k}=z^{k}+t^{k}\).
        \item Inner loop (\(N_{\mathrm{in}}\) steps): compute
        \(x=f_{\theta}(x_{\mathrm{LR}})\), \(y=\mathrm{Enc}(x)\),
        \(\hat{x}=\mathrm{Dec}(y)\), and update \(\theta\) by minimising
        \(\mathcal{L}\) in~\eqref{eq:training_loss}.
        \item Project: \(z^{k+1}=\Pi_{\mathcal{B}}(x-t^{k})\)
        via~\eqref{eq:projection_step}.
        \item Dual update: \(t^{k+1}=t^{k}+z^{k+1}-x\).
    \end{enumerate}
    \item Select the operating point, e.g.\ the iterate with lowest bitrate
    among those whose compressed PSNR stays within about \(1\,\mathrm{dB}\) of
    the quality budget.
\end{enumerate}

\subsection{Backbones and implementation}

The same loop covers four branches:
\begin{itemize}
    \item \textbf{NinaSR} at \(\times 2\) and \(\times 4\) (full fine-tuning);
    \item \textbf{Swin2SR} (full fine-tuning);
    \item \textbf{Swin2SR+LoRA} (frozen backbone; adapters on attention modules,
    e.g.\ rank \(8\), \(\alpha=16\), targets \texttt{query}, \texttt{value}).
\end{itemize}

In the reference setting the frozen codec is CompressAI
\texttt{cheng2020-anchor} at quality level \(6\)~\cite{cheng2020learned,begaint2020compressai}.
Typical grids use
\[
\lambda_1, \lambda_2 \in \{0.3, 0.5, 0.7, 1.0\}, \qquad
\mathrm{lr} \in \{10^{-4}, 3\cdot 10^{-4}, 10^{-3}\},
\]
with about \(N_{\mathrm{in}}=20\) inner and \(K=20\) outer steps per run on
screened DIV2K crops spanning a broad bpp range.
